\section{Theoretical Analysis}

\subsection{Computational Complexity}

\subsubsection{Per-Round Complexity}

The computational complexity of our Dual Scheduling mechanism per round is:

\textbf{Shapley Value Computation:}
\begin{itemize}
\item MC-Shapley with $M$ Monte Carlo iterations: $O(K \times M)$
\item For $K=10$ selected clients and $M=5$ iterations: $O(50)$
\item Significantly lower than exact Shapley: $O(2^K) = O(1024)$
\end{itemize}

\textbf{Energy and Channel Simulation:}
\begin{itemize}
\item Channel gain generation: $O(N)$
\item Energy consumption calculation: $O(K)$
\item Virtual queue update: $O(N)$
\end{itemize}

\textbf{Client Selection:}
\begin{itemize}
\item Score computation: $O(N)$
\item Top-K selection: $O(N \log K)$
\end{itemize}

\textbf{Total per-round complexity:} $O(N \log K + K \times M)$, which is practical for large-scale federated learning.

\subsection{Communication Complexity}

Each round requires:
\begin{itemize}
\item \textbf{Downlink}: Broadcasting global model $w_t$ to $K$ clients: $O(K \times d)$
\item \textbf{Uplink}: Receiving encrypted updates from $K$ clients: $O(K \times d)$
\item \textbf{Encryption overhead}: AES-GCM ciphertext expands plaintext by a fixed 28 bytes (12-byte nonce + 16-byte tag) per upload, negligible relative to gradient size $d \approx 270{,}000$ parameters
\item \textbf{Total}: $O(K \times d)$ per round, same asymptotic complexity as FedAvg
\end{itemize}

Our method does not materially increase communication overhead compared to baseline federated learning.

\subsection{Convergence Analysis}

Under standard assumptions (L-smooth loss functions and bounded gradient variance for non-IID data), our algorithm achieves convergence rate $O(1/\sqrt{T})$, identical to standard FedAvg. This follows from: (1) intelligent client selection reduces gradient variance compared to random sampling; (2) Shapley-based selection prioritizes high-contribution clients, accelerating convergence; (3) AES-256-GCM encryption operates purely at the transmission layer and does not perturb gradient values, introducing zero accuracy degradation compared to CDP-based approaches. Detailed proofs follow standard federated learning analysis techniques.
